\section{Threat model, limitations, and interpretation}
The theory and the evidence have different scopes. Proposition~1 and its factorization corollary hold for finite exact-terminal decision problems, but they identify when a target is attainable from a representation and output alphabet; they do not show that a deployed verifier possesses a target-sufficient representation. The empirical layers establish only the frozen representations, relations, and case distributions evaluated here.

First, the V2 and V3 campaigns use 420 synthetic mechanical-gold cases rather than naturally occurring scientific disagreements. The written human-adjudication rubric was never triggered because all final cases were mechanically decidable, so no claim is made about expert agreement. P4-X broadens the exact-contract battery to five artifact domains and eight hostile/control archetypes, but its 400 cases are still constructed contracts, not a population sample of scientific disputes. The final scoring logic is deterministic and uses no language-model provider; it isolates authority-transition semantics rather than the error distribution of a full stochastic research agent.

Second, the comparator mechanisms are matched reimplementations. They instantiate published ideas under one candidate-visible packet and common accounting, but the results are not a leaderboard against the original ProvenAI, ProvenanceGuard, DeepSciVerify, FIRE, AttributionBench, CLAIM-BENCH, RewardHackingAgents, or related systems in their native environments. Likewise, 60/60 clean coverage shows that the governed pipeline did not obtain its V2 false-promotion advantage by blanket refusal; the clean controls are intentionally unambiguous and do not estimate recall on difficult benign cases. The hostile taxonomy is finite, so adaptive attacks, compromised host infrastructure, unobserved side channels, and novel evaluator attacks remain outside the campaign.

Third, two registered V2 axes are saturated and must be read as design limits rather than comparative findings. Clean coverage is 1.0 for all eleven primary systems and every registered ablation. The original H3 family is also saturated: an empty evidence list makes every panel member correct on all 30 eligible cases. H1 therefore supplies the V2 comparative measurement, whereas H2's non-inferiority pass and V2 H3's null do not establish equality of underlying competences. The V3 construction repairs the registered H3 shortcuts and produces nonzero panel resolution, but nine of ten comparator mechanisms cannot emit the undetermined terminal. Its result therefore measures terminal/interface attainability under the frozen gate lattice. The $1.0$ margin against the H1-selected comparator and the $0.5$ margin against the escalation-capable mechanism are both retained precisely to prevent a broader judgement interpretation.

The nuisance audit has the same bounded scope. Fourteen frozen probes, evaluated across thirteen registered seeds, found no advantage on the exact V3 undetermined axis. That is evidence against the declared shortcut class, not a proof that no possible nuisance decoder exists. More generally, the theory separates semantic attainability, nuisance resistance, matched terminal attainability, and panel resolution because failure of any one of them changes what a benchmark contrast can mean.

Fourth, the component mechanisms are not claimed as novel. Provenance tracking, attribution evaluation, citation checking, iterative verification, evidence escalation, contamination detection, evaluator-tamper detection, abstention, benchmark auditing, claim-level auditability, provenance ontologies, and generic behavioral-integrity checking all have prior art discussed in Section~\ref{sec:related}. The contribution tested here is the non-compensatory composition of these prerequisites and, in P4-X, the target-bound scientific-promotion relation above a donor-complete generic authorization product. The information-equivalent typed product's 400/400 tie is a positive boundary result: the evidence does not support a centralization advantage, universal necessity of the exact gate list, or superiority over deployed donor systems.

Finally, naturalistic transfer remains unestablished. A prospective programme specifies 768 source--case-family clusters across four domains and eight information-state mechanisms, paired identifiable controls, worst-domain analysis, source-disjoint replication, and independent evaluator custody. The programme has so far produced only source-feasibility and identity-transport evidence. Its strongest current public identity bridge reaches 76/80 within one bounded software-publication frame, while author-lineage adjudication, same-claim natural-pair eligibility, source-disjoint replication, external custody, and exact external-comparator execution remain unresolved; zero natural pairs have been authorized. These preflights localize what a future naturalistic study would need, but they do not contribute performance evidence to the present paper. The detailed transport history is retained in the prospective supplementary record rather than used as part of the main decision proof.

Computational integrity is also not scientific authority. Frozen digests, deterministic replays, custody controls, and independent code paths establish that a stated computation was the computation performed and can be reconstructed under the released conditions. They do not by themselves establish that the benchmark measures the intended construct or that an external party has independently validated the scientific interpretation. A perfectly replayed flawed protocol would reproduce the flaw exactly; this is why the V2 H3 construction failure remains visible rather than being repaired retrospectively.
